\begin{pastexam}[label={px:ncu-114-dsa-20}]{中央資工所 114 學年度 第 20 題}
（複選；每題 5 分，全部答對才給分，答錯倒扣 1 分）

Choose correct statement(s):
\begin{enumerate}[label=(\alph*)]
\item Breadth-First Search (BFS) and Depth-First Search (DFS) are non-uninformed search strategies.
\item A* algorithm combines the actual cost from the start node ($g(n)$) and the heuristic estimate to
the goal ($h(n)$). If the heuristic is admissible (never overestimates the true cost) and consistent
(satisfies the triangle inequality), A* always guarantees finding the optimal solution.
\item Hill Climbing can escape from a local optimum due to its greedy nature.
\item Branch and Bound uses both upper and lower bounds to prune the search space for saving
computation.
\item Dijkstra's algorithm guarantees the shortest path even if the graph contains cycles.
\end{enumerate}
\end{pastexam}

\begin{pastsoln}
答案：(b)(d)(e)。

本題是搜尋策略的觀念。 BFS、DFS 在資結書第~21~章，分支限界在第~20.2~節，Dijkstra 在第~25.2~節；A* 和
爬山法（hill climbing）是人工智慧課的搜尋法，資結書沒有，這裡補充需要的部分。

\begin{itemize}
\item (a) 錯。 BFS、DFS 只照圖的結構展開，不用任何「離目標還多遠」的估計，是\textbf{無資訊}（uninformed，
又叫 blind）搜尋。 「non-uninformed」（有資訊）指的是 A*、最佳優先搜尋這類用啟發函數的方法。
\item (b) 對。 A* 每次展開 $f(n) = g(n) + h(n)$ 最小的節點。 $h$ 可容許（admissible，不高估）時，目標
第一次被展開時找到的就是最佳解；$h$ 再滿足一致性（consistent，$h(u) \le w(u, v) + h(v)$）時，沿任何
路徑 $f$ 不會變小，每個節點第一次被展開時 $g$ 就已經最小，用圖搜尋（不重複展開）也保證最佳。 （前提是
邊的成本為正、每個節點的分支數有限，這是 A* 的標準假設。） $h \equiv 0$ 時 A* 就是 Dijkstra。
\item (c) 錯。 爬山法每一步只移到比目前更好的鄰居，正因為這種貪心的性質，到了局部最佳（所有鄰居都
不比它好）就停住，\textbf{逃不出}局部最佳。 要逃出來得靠隨機重新開始、模擬退火這類額外的手段。
\item (d) 對。 分支限界（資結書第~20.2~節）以最小化問題為例：目前找到的最好解提供上界，每個節點算出
子樹裡所有解的下界；下界已經不比上界好的節點整棵剪掉。 兩種界一起用才能剪枝。
\item (e) 對。 Dijkstra 的正確性只要求邊的權重非負（資結書第~25.2~節），圖有沒有環無關；權重非負時
環不會讓路徑變短，最短路徑一定是簡單路徑。 （題目沒說權重，這裡照 Dijkstra 的前提，權重非負。）
\end{itemize}
\end{pastsoln}
